\section{实验}
\label{sec:experiments}

\subsection{实验设置}
\label{subsec:setup}

\paragraph{评估指标}
我们从图拓扑合法性、坐标回归精度和节点类型预测三个维度评估各子模块，
并以图结构分布指标衡量端到端生成质量。
所有指标均在 8,000 条独立测试集上计算。

\begin{itemize}
    \item \textbf{图有效率（Graph Validity Rate, GVR）}：
    生成图中连通且无孤立节点的样本占比，反映图拓扑生成的合法性。

    \item \textbf{度分布 KL 散度（Degree KL）}：
    生成图与真实图的节点度分布之间的 KL 散度，
    衡量端到端生成的图结构分布质量。

    \item \textbf{节点数分布 KL 散度（Node-count KL）}：
    生成图与真实图的节点数分布之间的 KL 散度，
    衡量模型对平面图规模的建模能力。

    \item \textbf{坐标 RMSE（Coord-RMSE）}：
    给定真实图结构条件下，去噪预测坐标与真实坐标的均方根误差（像素单位），
    衡量节点坐标扩散子模块的回归精度。

    \item \textbf{类型准确率（Type Acc.）}：
    给定真实坐标和图结构条件下，节点类型预测与真实类型的匹配率，
    衡量类型预测子模块的分类精度。
\end{itemize}

\paragraph{基线方法}
由于本文与现有方法在输入模态上存在根本差异，对比以方法特性说明为主：

\begin{itemize}
    \item \textbf{Architext}~\cite{galanos2023architext}：
    在 Architext\_v1 数据集上微调 GPT 系列模型，以自然语言描述直接生成平面图。
    其文本描述由规则模板生成，语言单一，且不建模显式图结构。

    \item \textbf{HouseDiffusion}~\cite{shabani2023housediffusion}：
    以气泡图（bubble diagram）为条件，通过混合扩散模型生成矢量多边形平面图。
    几何质量优秀，但需手工绘制气泡图，不支持自然语言输入。

    \item \textbf{ChatHouseDiffusion}~\cite{qin2025chathousediffusion}：
    在 HouseDiffusion 基础上引入 prompt 引导，支持文本条件生成与编辑，
    但缺乏对拓扑结构的显式约束。
\end{itemize}

\paragraph{实现细节}
\todo{待补充：训练硬件、训练时长、推理时间。}

\subsection{主实验结果}
\label{subsec:main_results}

\todo{待写：与 Architext 基线的对比。}

\subsection{消融实验}
\label{subsec:ablation}

\todo{待写。}

\subsection{定性结果}
\label{subsec:qualitative}

\todo{待写：生成结果可视化。}
